\section{Joint Physical and Measurement Model}
\label{sect:model}

\subsection{Physical Parameterization}

The fitted binary parameter vector is
\begin{equation}
\Theta=\{P,T_0,e,i,\Omega,\omega_1,M_1,M_2,\varpi,\gamma,\betag\},
\label{eq:theta}
\end{equation}
where $\omega_1$ is the primary spectroscopic argument of periastron and $\Omega$ is measured from local North through East. The relative semi-major axis follows from Kepler's third law,
\begin{equation}
a_{\rm rel}^3=\frac{G(M_1+M_2)P^2}{4\pi^2},
\end{equation}
and the barycentric component axes are
\begin{equation}
a_1=a_{\rm rel}\frac{M_2}{M_1+M_2},\qquad
a_2=a_{\rm rel}\frac{M_1}{M_1+M_2}.
\end{equation}
The RV semi-amplitudes are derived rather than fitted independently,
\begin{equation}
K_j=\frac{2\pi a_j\sin i}{P\sqrt{1-e^2}}.
\end{equation}

At time $t$, the mean anomaly is
\begin{equation}
M(t)=2\pi\frac{t-T_0}{P},
\end{equation}
the eccentric anomaly satisfies $M=E-e\sin E$, and
\begin{equation}
\nu=2\arctan2\!\left(\sqrt{1+e}\sin\frac{E}{2},
\sqrt{1-e}\cos\frac{E}{2}\right).
\end{equation}

\subsection{Orbit Conventions and Relative Astrometry}

A joint visual/SB2 model is vulnerable to a $180^\circ$ convention error if the same numerical argument of periastron is used simultaneously for the primary RV curve and the relative vector $\mathbf r_2-\mathbf r_1$. We adopt
\begin{equation}
\omega_{\rm rel}=\omega_1+\pi,
\label{eq:omegarel}
\end{equation}
consistent with the primary-star Campbell convention used in combined orbit work \citep{Leclerc2023}. For $u=\nu+\omega_{\rm rel}$ and instantaneous separation $r$, the relative North and East coordinates are
\begin{align}
\Delta N &= r(\cos\Omega\cos u-\sin\Omega\sin u\cos i),\\
\Delta E &= r(\sin\Omega\cos u+\cos\Omega\sin u\cos i),
\end{align}
and the tangent-plane order used by the code is $(\Delta\alpha^\ast,\Delta\delta)=(\Delta E,\Delta N)$. The line-of-sight coordinate is positive away from the observer and is regression-tested to satisfy
\begin{equation}
\frac{dz_{\rm rel}}{dt}=RV_2-RV_1.
\end{equation}

For resolved astrometry, the angular model is obtained by scaling the physical relative vector with the parallax. The present synthetic experiments use isotropic $2\times2$ astrometric covariance matrices, while the likelihood implementation supports full covariance whitening.

\subsection{SB2 Radial Velocities}

The spectroscopic factor is
\begin{equation}
F(t)=\cos[\nu(t)+\omega_1]+e\cos\omega_1,
\end{equation}
and the component velocities are
\begin{equation}
RV_1=\gamma+K_1F(t),\qquad RV_2=\gamma-K_2F(t).
\label{eq:rv}
\end{equation}
This convention satisfies barycentric momentum balance,
\begin{equation}
M_1(RV_1-\gamma)+M_2(RV_2-\gamma)=0.
\end{equation}

\subsection{Photocentre Model $\mzero$}

Define the secondary mass fraction and Gaia-band light fraction as
\begin{equation}
B=\frac{M_2}{M_1+M_2},\qquad
\betag=\frac{F_{2,G}}{F_{1,G}+F_{2,G}}.
\end{equation}
For the relative vector $\mathbf r=\mathbf r_2-\mathbf r_1$, the unresolved photocentre relative to the barycentre is
\begin{equation}
\mathbf r_{\rm ph}=(\betag-B)\mathbf r.
\label{eq:photocentre}
\end{equation}
For directional Gaia scan angle $\psi$, defined with $0^\circ$ toward North and $90^\circ$ toward East, the projected relative separation is
\begin{equation}
\Delta\eta=\Delta\alpha^\ast\sin\psi+\Delta\delta\cos\psi.
\label{eq:alproj}
\end{equation}
Model $\mzero$ therefore assigns the orbital along-scan displacement $(\betag-B)\Delta\eta$.

\subsection{Equal-Width Blended Response $\mone$}

The first resolution-aware model represents the one-dimensional light profile as
\begin{equation}
I(x)=F_1\exp\!\left[-\frac{(x-x_1)^2}{2\alpha^2}\right]
+F_2\exp\!\left[-\frac{(x-x_2)^2}{2\alpha^2}\right],
\label{eq:gaussian}
\end{equation}
and takes the measured coordinate to be the unique global maximum of $I(x)$. This is the same equal-width response family implemented in \texttt{gaiamock} \citep{ElBadry2024}; its use here is an implementation of published measurement physics, not a novelty claim. In the original baseline experiments the width was denoted $\sigma$. For the response-fidelity experiments we write the narrow/along-axis research width as $\alpha$ so that it can be distinguished from the finite across-axis width introduced below.

The equal-width mixture has an exact single-to-double-mode boundary. For $F_1=1-\betag$, $F_2=\betag$ and projected separation $s$, the critical ratio follows from simultaneous first- and second-derivative conditions:
\begin{equation}
\sinh(2u)-2u=\ln\!\left(\frac{1-\betag}{\betag}\right),
\end{equation}
\begin{equation}
\frac{s_{\rm crit}}{\alpha}=2\cosh u.
\label{eq:scrit}
\end{equation}
Equation~(\ref{eq:scrit}) is used only as a validity criterion for this restricted equal-width profile. It is not a calibrated Gaia angular-resolution threshold.

\subsection{Orientation-Dependent Response $\mtwo$}

To test response fidelity rather than matched-model recovery alone, the injected model is generalized using the elongated-Gaussian treatment of \citet{Penoyre2026}. Let $\alpha$ denote the narrow/along-axis width and $\beta_{\rm PSF}\geq\alpha$ the long/across-axis width of the idealized Gaussian. If $\phi$ is the angle between the binary separation vector and the scan axis, the effective one-dimensional width is
\begin{equation}
\gamma(\phi)=
\left[
\frac{\cos^2\phi}{\alpha^2}
+\frac{\sin^2\phi}{\beta_{\rm PSF}^2}
\right]^{-1/2}.
\label{eq:gamma}
\end{equation}
The two-dimensional blend can then be reduced to the same dimensionless one-dimensional peak problem after scaling the full component separation $r$ by $\gamma(\phi)$. The code solves the peak numerically using the exact effective width in Equation~(\ref{eq:gamma}); it does not use the corrected low-order series expressions from \citet{Penoyre2026}.

For finite $\beta_{\rm PSF}$, the mode-count criterion uses
\begin{equation}
\rho=\frac{r}{\gamma(\phi)}
\end{equation}
together with the same light-ratio-dependent critical value as the equal-width one-dimensional mixture. The resulting peak location along the pair axis is projected back onto the along-scan direction. In the limit
\begin{equation}
\beta_{\rm PSF}\rightarrow\infty,
\end{equation}
the implementation is regression-tested to recover $\mone$ exactly. This is the appropriate limiting connection to the effectively unconstrained across-scan assumption of the one-dimensional response; a circular Gaussian with $\beta_{\rm PSF}=\alpha$ is not identified with $\mone$.

Neither $\alpha$ nor $\beta_{\rm PSF}$ is a calibrated Gaia PLSF width. The finite-elongation model remains an idealized research surrogate. Gaia's actual PLSF contains additional time-, colour-, focal-plane-, window-, and drift-scan dependences \citep{Rowell2021,Rowell2026}.

\subsection{Single-Peak Validity}

A one-coordinate likelihood becomes ambiguous when the surrogate profile is multi-peaked. Before each response-fidelity realization, epochs classified as multi-peak under the injected $\mtwo$ response are excluded. All fitted models then use the identical retained epochs. The optimizer may traverse a continuous peak branch internally, but an $\mone$ or $\mtwo$ fit is accepted for scientific analysis only when the final fitted response is single-peaked at every retained epoch. The reported 720-fit, 270-fit, and 100-seed experiments all retain all 87 Gaia-like epochs and contain no multi-peak injection or final prediction.

\subsection{Joint Objective and Paired Model Comparison}

The deterministic validation objective combines resolved relative astrometry, both SB2 velocities, and Gaia-like along-scan observations. All eleven parameters in Equation~(\ref{eq:theta}) are fitted with bounded least squares. For each physical configuration and random seed, $\mzero$, $\mone$, and $\mtwo$ are fitted to the same synthetic realization.

We use
\begin{align}
\Delta\chi^2_{0,2}&=\chi^2_{\mzero}-\chi^2_{\mtwo},\\
\Delta\chi^2_{1,2}&=\chi^2_{\mone}-\chi^2_{\mtwo}
\end{align}
as descriptive paired mismatch diagnostics. Positive values favor the more faithful $\mtwo$ response on the same data, but the models are not treated as nested and the differences are not converted to Gaussian detection significances.

The primary scientific diagnostics are instead the fractional parameter biases,
\begin{equation}
\delta_x=\frac{\widehat{x}-x_{\rm true}}{x_{\rm true}},
\label{eq:fracbias}
\end{equation}
especially for the physical component masses $M_1$ and $M_2$, parallax $\varpi$, and light fraction $\betag$. The current study uses deterministic fits to isolate response-model bias. Posterior coverage is intentionally deferred until this hierarchy is established.
